\problema{Meeting Rooms III}{2402}{src/03-intervalos/2402-meeting-rooms-iii.cpp}{H}

Hay \texttt{n} salas numeradas de \texttt{0} a \texttt{n-1}. Nos dan
\texttt{meetings[i] = [inicio, fin]}, y se procesan en orden de
\texttt{inicio}. Cada reunión se asigna a la sala \emph{libre} de menor
índice disponible en ese momento. Si ninguna está libre, la reunión se
\emph{retrasa} hasta que la sala que se desocupe primero quede disponible,
conservando la \emph{misma duración} que tenía originalmente, y se asigna a
esa sala (la de menor índice si varias se desocupan al mismo tiempo).
Queremos la sala que organizó más reuniones en total (la de menor índice en
caso de empate).

\paragraph{La idea, con un ejemplo de la vida real.} Es la versión
``estricta'' de repartir salas: en vez de simplemente contar cuántas salas
hacen falta (como en el problema anterior), aquí ya existe un número fijo
de salas y hay que decidir, reunión por reunión, cuál usa cada una,
incluyendo qué pasa cuando llega una reunión y no hay ninguna sala libre en
ese instante: se le dice a los organizadores ``espera a que alguna sala se
desocupe'', y su reunión se recorre en el tiempo sin acortarse. Es como una
recepcionista con \texttt{n} salas de juntas que siempre prefiere la sala 1
sobre la sala 2 cuando ambas están libres, y que sabe exactamente cuál sala
va a desocuparse primero cuando todas están ocupadas.

\noindent Con \texttt{n = 2}, \texttt{meetings = [[0,10],[1,5],[2,7],[3,4]]}:
\texttt{[0,10]} toma la sala 0, \texttt{[1,5]} toma la sala 1. Cuando llega
\texttt{[2,7]} ninguna sala está libre; la sala 1 se desocupa primero (en el
tiempo 5), así que la reunión se retrasa: dura $7-2=5$, ocupa la sala 1 de
5 a 10. Cuando llega \texttt{[3,4]} tampoco hay sala libre; ambas se
desocupan en el tiempo 10 (empate), gana la de menor índice, sala 0; dura
$4-3=1$, ocupa la sala 0 de 10 a 11. Cada sala terminó con 2 reuniones:
empate, gana la sala 0.

\paragraph{Cómo lo resuelve el código, paso a paso.}
\begin{enumerate}
  \item Ordenamos las reuniones por \texttt{inicio}, porque el enunciado
        exige procesarlas en ese orden.
  \item Mantenemos dos \emph{min-heaps} (colas de prioridad que siempre dan
        el menor elemento primero): uno de \textbf{salas libres}, ordenado
        por índice de sala; otro de \textbf{salas ocupadas}, ordenado por
        la hora en que cada una se desocupa (y, en empate, por índice).
  \item Antes de procesar cada reunión, sacamos del heap de ocupadas todas
        las salas cuya hora de fin ya pasó (es decir, es menor o igual al
        inicio de la reunión actual) y las movemos al heap de libres.
  \item Si después de eso hay alguna sala libre, la reunión la usa (la de
        menor índice, que es justo lo que da un min-heap por índice) y esa
        sala vuelve al heap de ocupadas con su hora de fin real.
  \item Si no hay ninguna sala libre, tomamos del heap de ocupadas la que se
        va a desocupar \emph{más pronto}: la reunión se retrasa hasta esa
        hora, conservando su duración original, y esa sala vuelve al heap
        de ocupadas con la nueva hora de fin (hora en que se liberó más la
        duración).
  \item En cualquiera de los dos casos, sumamos uno al contador de esa sala.
  \item Al final, devolvemos la sala con mayor contador (la de menor índice
        en caso de empate).
\end{enumerate}

\lstinputlisting{src/03-intervalos/2402-meeting-rooms-iii.cpp}

\complejidad{$O(n \log n + n \log k)$}{$O(n + k)$}

\paragraph{¿Qué significa esa complejidad?} Ordenar las $n$ reuniones cuesta
$O(n \log n)$. Cada reunión hace un número constante de operaciones sobre
heaps de tamaño a lo más $k$ (el número de salas), así que procesarlas
todas cuesta $O(n \log k)$. En espacio guardamos los dos heaps (tamaño total
acotado por $k$) más el arreglo de contadores, tamaño $k$, además de una
copia ordenada de las reuniones, tamaño $n$.

\paragraph{Consejos para la entrevista (Amazon).} Este problema mezcla dos
ideas del bloque de intervalos: el barrido por tiempo (como en Meeting
Rooms II) y el uso de dos heaps con roles distintos (como en Find Median
from Data Stream). El punto más fácil de arruinar es la regla de retraso:
hay que conservar la \emph{duración}, no la hora de fin original, al mover
una reunión en el tiempo. Otro detalle fino: cuando varias salas se
desocupan exactamente a la misma hora, el desempate por índice debe
respetarse tanto al elegir sala libre como al elegir cuál sala ocupada se
desocupa primero; por eso el heap de ocupadas ordena por el par
\texttt{(hora de fin, índice)} y no solo por la hora. Menciona el caso
límite de una sola sala, donde todas las reuniones, se traslapen o no,
terminan cayendo ahí, posiblemente retrasadas en cadena.
